\providecommand{\stepbox}[1]{\parbox[t][2.75cm][t]{6.3cm}{\raggedright#1}}%
\begin{tikzpicture}[sb, x=1cm, y=1cm,
  step/.style={#1, font=\small, inner sep=5pt},
  bundle/.style={font=\footnotesize, text=sbOther, align=center,
    text width=7.2cm, align=flush center},
  uses/.style={draw=sbInk, fill=white, rounded corners=1.5pt,
    font=\scriptsize, inner sep=2pt},
  usesSN/.style={uses, fill=sbBound!40},
  flow/.style={->, sbInk, line width=1.0pt},
  lab/.style={font=\footnotesize, fill=white, inner sep=1.5pt, align=center},
]
\node[step=sb aux box] (s1) at (0,0) {\stepbox{%
  \textbf{Step 1: Siegel lemma.}\\[2pt]
  $u_1,\dots,u_r$ an adapted basis of $\calL(aP_0)$,\\
  $\Lambda\colon\calL(bP_0)^r\to\calL\bigl(MaP_0+b\,\sigma(P_0)\bigr)$,\\
  $\Lambda(v_i)=\sum_iu_i^M\sigma(v_i)$.\\[2pt]
  Source larger than target by (S), so some $(v_i)\ne0$ has
  $\Lambda(v_i)=0$.}};
\node[step=sb aux box] (s2) at (7.6,0) {\stepbox{%
  \textbf{Step 2: non-vanishing.}\\[2pt]
  $w=\sum_iu_iv_i^{M'}\in\calL\bigl((a+M'b)P_0\bigr)$.\\[2pt]
  By (N), the pole orders $-v_{P_0}(u_i)$ are distinct integers in
  $[0,M')$, so $w\ne0$
  (Proposition~\ref{prop:independence}).}};
\node[step=sb aux box] (s3) at (7.6,-5.0) {\stepbox{%
  \textbf{Step 3: forced multiplicity.}\\[2pt]
  $w^M=\sum_iu_i^Mv_i^q$, as $z\mapsto z^M$ is additive.\\[2pt]
  At a $\sigma$-Frobenius place $P\ne P_0$:
  $v_i^q\equiv\sigma(v_i)\bmod P$, so
  $w^M\equiv\Lambda(v_i)=0\bmod P$, hence $w\in P$.}};
\node[step=sb bound box] (s4) at (0,-5.0) {\stepbox{%
  \textbf{Step 4: slope inequality.}\\[2pt]
  $Z_\sigma=\sum P$ over the $\sigma$-Frobenius places $P\ne P_0$.\\[2pt]
  $(w)+(a+M'b)P_0-Z_\sigma\ge0$,\\ so $\deg Z_\sigma\le a+M'b$ and\\[2pt]
  \hfil$N(K,\sigma)\le1+a+M'b$.}};
\node[usesSN, anchor=south east] at ($(s1.south east)+(-3pt,3pt)$)
  {uses (S)};
\node[usesSN, anchor=south east] at ($(s2.south east)+(-3pt,3pt)$)
  {uses (N)};
\node[uses, anchor=south east] at ($(s3.south east)+(-3pt,3pt)$)
  {uses char $p$};
\node[uses, anchor=south east] at ($(s4.south east)+(-3pt,3pt)$)
  {uses \ref{S2}};
\node[bundle, below=2pt of s1] (b1)
  {$H^0(\calO_C(bP_0))^r\to H^0\bigl(\calO_C(MaP_0+b\,\sigma(P_0))\bigr)$};
\node[bundle, below=2pt of s2] (b2)
  {$w\in H^0\bigl(\calO_C((a+M'b)P_0)\bigr)$};
\node[bundle, below=2pt of s3] (b3)
  {$w^M\in H^0\bigl(\calO_C((Ma+qb)P_0)\bigr)$,\\ vanishing to order $\ge M$ at each $P$};
\node[bundle, below=2pt of s4] (b4)
  {$w\colon\calO_C(Z_\sigma)\to\calO_C\bigl((a+M'b)P_0\bigr)$, nonzero};
\draw[flow] (s1.east) -- (s2.west) node[lab, midway, above=1pt, fill=none]
  {$(v_i)$};
\draw[flow] (b2.south) -- (s3.north) node[lab, midway, right=2pt, fill=none]
  {$w$};
\draw[flow] (s3.west) -- (s4.east) node[lab, midway, above=1pt, fill=none,
  text width=1.1cm] {zeros at all $P$};
\draw[->, sbInk, line width=0.6pt, dashed]
  ($(b1.south east)+(-0.6,0)$) to[out=-40, in=150]
  node[lab, pos=0.45, above right=-1pt, fill=none] {$\Lambda(v_i)=0$}
  ($(s3.north west)+(0.6,0)$);
\end{tikzpicture}
\caption{The four steps of the proof of Theorem~\ref{thm:aux-section}, for a
function field $K/\F_q$, a place $P_0$ of degree one,
$\sigma\in\Aut(K/k)$ and $q=MM'$ with $M$ a power of $p$. Step~1 uses the
Siegel condition (S) to find a nonzero kernel element $(v_i)$ of
$\Lambda$; Step~2 uses the non-vanishing condition (N) and
Proposition~\ref{prop:independence} to see that $w\ne0$; Step~3 uses that
$z\mapsto z^M$ is additive, the only use of the characteristic
(Proposition~\ref{prop:frobenius}), and the relation $\Lambda(v_i)=0$ of
Step~1 (dashed arrow), to make $w$ vanish at every $\sigma$-Frobenius place
$P\ne P_0$; Step~4 compares degrees (fact~\ref{S2}). The tag in each box
names the hypothesis that step uses. Grey lines give the line bundles on
$C$ in which each step takes place. Sky blue: auxiliary data; orange: the
bound produced and the conditions (S), (N). Schematic.}
\label{fig:four-steps}
